\section{企业服务级 AI：公开 token 会卷，私有数据有价值}

前两章解释了明略为什么拥有数据入口、交付经验和经营耐力，本章进入整期的技术重心，问题变成：这些积累在 AI 时代如何转成企业级 Agentic Model。吴明辉判断，基于公开数据训练基础模型、以卖 token 为商业模式的公司会非常卷，最后容易卷成电费。因为公开数据和通用模型能力会趋同，token 价格会下降。企业级 AI 的差异化不在于重复训练一个公开模型，而在于私有 data、行业流程和真实工具系统。

\begin{figure}[H]
\centering
\includegraphics[width=0.94\textwidth]{private-data-value.png}
\caption{私有 Data 的差异化价值：公开数据卖 token 会卷，私有数据能形成行业价值。自制概念图，依据 01:45:01--02:06:40 对谈内容整理。}
\end{figure}

\begin{knowledgebox}{读图：为什么私有数据能产生溢价}
公开数据训练出的能力会越来越商品化；企业私有数据包含客户、流程、交易、合同、对话、设备和风险记录。模型只有进入这些上下文，才可能生成企业愿意付费的差异化结果。
\end{knowledgebox}

\begin{warningbox}{“私有数据”不是简单上传企业文件}
企业私有数据真正难的部分在权限、实时性、结构化、脏数据、业务含义和责任边界。把文档丢进向量库只是第一步；让 Agent 正确理解谁能看、谁能改、什么结果可交付，才是企业级 AI 的门槛。
\end{warningbox}

\begin{knowledgebox}{企业 AI 的四层架构}
\begin{tabularx}{\textwidth}{p{0.20\textwidth}p{0.32\textwidth}X}
\toprule
层级 & 作用 & 典型失败方式 \\
\midrule
数据层 & 接入私有数据、权限和历史记录 & 数据脏、权限错、上下文不完整。 \\
模型层 & 理解、推理、生成候选动作 & 只会总结，不懂业务约束。 \\
工具层 & 调用业务系统、数据库和流程接口 & 无法写回系统，停留在建议。 \\
治理层 & 审计、复核、责任和反馈 & 企业不敢放权，无法规模化部署。 \\
\bottomrule
\end{tabularx}
\end{knowledgebox}

\begin{knowledgebox}{从通用模型到企业交付的转译链}
\begin{tabularx}{\textwidth}{p{0.22\textwidth}p{0.32\textwidth}X}
\toprule
转译环节 & 要回答的问题 & 明略式经验的价值 \\
\midrule
行业语义 & 企业说的“客户”“风险”“线索”到底指什么？ & 长期 ToB 交付积累了行业语义和客户现场理解。 \\
流程状态 & 任务处在售前、履约、回款还是售后？ & CRM、销售、客服对话数据提供流程轨迹。 \\
结果计量 & 动作是否真的改善收入、成本或风险？ & 秒针和数据分析经验强调可计量、可审计。 \\
组织采纳 & 谁批准、谁复核、谁承担后果？ & ToB 经营训练了权限、责任和交付意识。 \\
\bottomrule
\end{tabularx}
\end{knowledgebox}

这里还隐含了一个商业模式判断：如果模型公司只卖公开 token，价格会被供给增加和推理优化不断压低；如果企业服务公司把 token 变成可计价的业务动作，单价就不再只由算力成本决定。例如，同样消耗一百万 token，生成一段普通总结和发现一个重大合同风险的商业价值完全不同。企业级 AI 的定价单位最终可能从 token 转向任务、风险、收入提升和成本节约。

\begin{importantbox}{从 token 定价到结果定价}
公开模型 API 的计价单位通常是 token；企业级 AI 的理想计价单位应该是结果。只有当模型输出能进入 workflow、产生可验证业务收益，企业才会为“结果”而不是“字数”付费。
\end{importantbox}

\subsection{现实世界的数值游戏}

上一节讲私有数据，本节解释为什么企业数据难以直接转成模型能力。吴明辉把企业服务称为现实世界的数值游戏。围棋很适合 AI，是因为棋盘透明、规则完备、context 全；企业工作则相反：上下文不完整、信息分散、目标模糊、流程跨系统、结果要对真实客户和财务负责。企业级 AI 的难点不是“会不会回答”，而是能不能把不完整上下文转成可执行动作。

\begin{importantbox}{企业级 AI 的难点公式}
\[
\text{Enterprise AI Value} \approx f(\text{Private Data}, \text{Workflow}, \text{Tool Use}, \text{Trust})
\]
其中，\(\text{Private Data}\) 是企业私有上下文，\(\text{Workflow}\) 是业务流程，\(\text{Tool Use}\) 是调用系统和工具的能力，\(\text{Trust}\) 是可审计、可控、可交付的信任。
\end{importantbox}

\begin{knowledgebox}{术语消化：企业级 AI 的四个约束}
\begin{tabularx}{\textwidth}{p{0.20\textwidth}p{0.32\textwidth}X}
\toprule
约束 & 含义 & 产品后果 \\
\midrule
Private Data & 企业独有的客户、流程、交易和对话数据 & 决定模型是否懂业务现场。 \\
Workflow & 跨部门、跨系统的真实流程 & 决定 Agent 是否能从建议走向交付。 \\
Tool Use & 调用 CRM、ERP、工单、数据库等工具 & 决定 AI 是否能改变业务状态。 \\
Trust & 权限、审计、复核和责任边界 & 决定企业是否敢把任务交给 AI。 \\
\bottomrule
\end{tabularx}
\end{knowledgebox}

这也是为什么“现实世界的数值游戏”比围棋更难。围棋的状态空间很大，但棋盘完整、规则明确、胜负可判；企业业务的状态空间不仅大，而且很多变量藏在人的脑子里、客户关系里、历史合同里和跨系统流程里。Agent 要想在这里工作，必须先把隐性知识转成可被系统利用的 memory，再把 memory 与工具调用结合起来。

\subsection{Agentic Model：从模型能力到工作流执行}

本节解释 Agentic Model。访谈里吴明辉把 Agent 或 Tool Use 看成企业服务领域的新链接。模型不再只是回答问题，而是接入数据、理解任务、调用工具、完成流程并返回可审计结果。这和 EP139 的 Agent loop、EP118 的 Agent OS 形成呼应，但企业场景更强调私有化、权限、审计和责任边界。

\begin{figure}[H]
\centering
\includegraphics[width=0.96\textwidth]{enterprise-agentic-model.png}
\caption{企业级 Agentic Model：从私有数据到工具调用，再到企业工作流执行。自制概念图，依据 01:45:01--03:00:00 对谈内容整理。}
\end{figure}

\begin{knowledgebox}{读图：企业 Agent 不是聊天机器人}
从左到右看，企业 Agent 先要拿到私有数据和企业上下文，再用模型理解和推理，再通过 Tool Use 调用 CRM、ERP、工单、数据库等系统，最后进入可审计工作流。缺一环都很难成为企业级产品。
\end{knowledgebox}

\subsection{DeepMiner：把企业数据挖成行动线索}

Agentic Model 解决“怎样执行”，本节用 DeepMiner 说明“从哪里发现任务”。DeepMiner 在访谈中出现不多，但它是理解明略新产品方向的线索。它代表一种企业 AI 产品形态：不是只做问答，而是把多系统数据接入、挖掘规则/风险/机会、形成洞察，再通过 Agent 或工具执行后续动作。

\begin{figure}[H]
\centering
\includegraphics[width=0.96\textwidth]{deepminer-product-loop.png}
\caption{DeepMiner 产品闭环：用 AI 挖掘企业数据中的规则、风险和行动线索。自制概念图，依据 02:31:58--03:36:59 对谈内容整理。}
\end{figure}

\begin{knowledgebox}{读图：从数据挖掘到 Agentic 交付}
传统 BI/数据分析更偏洞察展示；DeepMiner 式产品如果接入 Agent，就可以从“发现问题”继续走到“提出行动”“调用系统”“记录结果”。这正是企业 AI 从分析工具变成生产系统的关键。
\end{knowledgebox}

\begin{importantbox}{DeepMiner 式产品的关键转变}
过去的数据分析常停在 dashboard 和报表；企业级 Agentic 产品必须进一步回答：发现了什么异常，应该采取什么动作，动作由谁批准，调用哪个系统，结果如何回流。
\end{importantbox}

\begin{knowledgebox}{从 Copilot 到 Agentic System}
\begin{tabularx}{\textwidth}{p{0.22\textwidth}p{0.34\textwidth}X}
\toprule
阶段 & 典型形态 & 企业侧要求 \\
\midrule
Copilot & 人问、AI 答，人复制结果 & 低风险、高频辅助，主要提升个人效率。 \\
Workflow Agent & AI 读数据、填表、发起流程 & 要接入权限、工具和审批。 \\
Agentic System & 多 Agent 协同完成任务 & 要有日志、调度、失败恢复和结果责任。 \\
Trusted Agentic Model & 可信地承担企业关键任务 & 要可审计、可复核、可私有部署或受控运行。 \\
\bottomrule
\end{tabularx}
\end{knowledgebox}

\subsection{本章小结}

企业级 AI 的核心在于私有数据和工作流。公开 token 会变便宜，真正有溢价的是能把企业上下文变成结果的 Agentic Model。

